\documentclass[12pt,a4paper]{article}

% ------------------------------- packages ----------------------------------
\usepackage[utf8]{inputenc}
\usepackage[T1]{fontenc}
\usepackage{lmodern}
\usepackage[margin=1in]{geometry}
\usepackage{amsmath,amssymb,amsthm}
\usepackage{mathtools}
\usepackage{microtype}
\usepackage{natbib}
\bibliographystyle{plainnat}
\usepackage[most]{tcolorbox}
\tcbuselibrary{breakable,skins}
\usepackage{hyperref}
\usepackage[nameinlink,capitalize]{cleveref}
\usepackage{enumitem}

% ------------------------------- theorems ----------------------------------
\theoremstyle{plain}
\newtheorem{theorem}{Theorem}[section]
\newtheorem{lemma}[theorem]{Lemma}
\newtheorem{corollary}[theorem]{Corollary}

\theoremstyle{definition}
\newtheorem{definition}[theorem]{Definition}
\newtheorem{example}[theorem]{Example}

\crefname{theorem}{Theorem}{Theorems}
\crefname{lemma}{Lemma}{Lemmas}
\crefname{corollary}{Corollary}{Corollaries}
\crefname{definition}{Definition}{Definitions}
\crefname{example}{Example}{Examples}
\crefname{section}{Section}{Sections}
\crefname{equation}{Equation}{Equations}

% ------------------------------- intuition box -----------------------------
\newtcolorbox{intuition}[1][]{
  breakable,
  enhanced,
  colback=yellow!8,
  colframe=orange!60!black,
  coltitle=white,
  fonttitle=\bfseries,
  title={Intuition},
  boxrule=0.6pt,
  left=6pt,right=6pt,top=4pt,bottom=4pt,
  #1
}

% ------------------------------- title ------------------------------------
\title{Prediction, Compression, and Intelligence:\\
       Three Names for the Same Operation}
\author{Section 3 --- Philosophy of Machine Minds\\
        Executive Education Lecture Notes}
\date{}

\begin{document}
\maketitle

% =========================================================================
\section*{Preface: how to read these notes}
% =========================================================================

These notes are written on two tracks at once, for two readers travelling
together.  The \emph{student} --- typically a bachelor-STEM reader who
wants the theorems, proofs, and definitions --- will find the main text
body rigorous and self-contained.  The \emph{manager} --- a
decision-maker who wants the conceptual story without the symbols ---
will find every formal object wrapped in a shaded \textbf{Intuition}
box written in plain English, with analogies, small numerical examples,
and a ``why should a decision-maker care'' coda.  Each Intuition box
cross-references the nearest definition or theorem via
\cref{sec:why}\ldots so a curious manager can always drop down into
the math, and a rigorous student can always drop up into the metaphor.
Read strictly down the page and you get both tracks interleaved; read
only the Intuition boxes and the section headings and you will still
leave with the thesis: \emph{prediction, compression, and intelligence
are, under precise definitions, the same operation, and that fact is
the implicit business plan of every major foundation-model lab.}

\tableofcontents
\newpage

% =========================================================================
\section{Why this equivalence matters}
\label{sec:why}
% =========================================================================

\begin{intuition}
The claim we unpack in these notes, in one sentence: \emph{being good
at predicting what comes next, being good at compressing a file, and
being good at thinking are --- up to translation --- the same skill.}
That sounds like a philosopher's trick.  It is not.  It is a chain of
three mathematical results, each modest on its own, that lock together
like gears.  Shannon showed prediction and compression are literally
the same job in different units (\cref{thm:ce-code}).  Kolmogorov
showed that ``the shortest description of a thing'' is a deep,
coordinate-free notion of what the thing \emph{is}
(\cref{def:kolmogorov,thm:invariance}).  Solomonoff and Hutter showed
that an agent who bets according to ``shorter explanations are more
likely'' is, in a precise sense, the smartest possible learner
(\cref{thm:solomonoff,thm:aixi-opt}).

Why should a decision-maker care?  Because your vendors are selling
you products whose entire training objective is one single number:
how surprised the model was by the next word of the internet.  If you
believe the chain above, minimising that one number is not a narrow
engineering goal --- it is a proxy for general intelligence.  That is
either the most important business fact of the decade or a confusion
between two different things sharing a name.  You cannot tell which
without understanding the equivalence.
Think of it as three currencies --- \textbf{bits saved} (how tightly
a model compresses your data), \textbf{dollars earned} (what that
compression is worth when deployed), and \textbf{capability on
held-out tasks} (how well the model handles problems it has not
seen) --- and three exchange rates between them.  Three currencies,
three exchange rates, one underlying quantity.  Once you see the
rates, a lot of hype becomes legible, and a lot of legitimate
excitement stops sounding like magic.

\textbf{One frame to carry throughout.}  A perfect next-token
predictor would be an ideal learner; real next-token predictors are
not perfect, and the gap is exactly where this document's caveats
live (\cref{sec:caveats}).  We state the thesis before we qualify
it, but the qualification is not optional.
\end{intuition}

The central technical claim of this note is that three operations that
are often treated as distinct --- predicting the next symbol in a
sequence, compressing that sequence into a shorter description, and
behaving intelligently in its presence --- are, under a precise set of
definitions, the \emph{same operation} viewed from three angles.
Shannon's source coding theorem \citep{shannon1948mathematical}
identifies prediction with compression; Kolmogorov complexity
\citep{kolmogorov1965three} makes the notion of ``shortest
description'' independent of the choice of language; Solomonoff's
universal prior \citep{solomonoff1964formal} promotes the Kolmogorov
minimiser to a Bayesian predictor that converges to any computable
source; and Hutter's AIXI \citep{hutter2005universal} upgrades that
predictor into a reward-seeking agent whose optimality can be stated
as a theorem.  The chain is not an analogy.  Every step in this note
is a definition, a lemma, or a theorem.  The pedagogical payoff is
that one can then ask sharp questions about large language models ---
whether they count as compressors, whether better compression entails
more general capability, and where the equivalence breaks --- without
smuggling in metaphor.

% =========================================================================
\section{Prelude: Shannon and the cost of surprise}
\label{sec:shannon-prelude}
% =========================================================================

\begin{intuition}
Information is the cost of being surprised, and you are surprised
exactly as much as an event was unlikely.  In Los Angeles, a
forecaster who says ``sunny'' every day will be right most of the
time and needs almost no words.  In Zurich, she needs rain, snow,
F\"{o}hn, fog, and sun.  The \emph{average} number of yes/no questions
she needs to convey a day's weather is the \emph{entropy}
(\cref{def:entropy}), denoted $H(X)$.  Flip a fair coin: one yes/no
question suffices, so $H = 1$~bit.  Flip a coin biased 7/8 heads:
most of the time you guess heads and are done; on average you need
only about half a bit per flip.  A 1-in-1024 event carries ten bits
of ``wow''; a 50/50 event carries one bit of ``huh.''

Shannon's source-coding theorem (\cref{thm:source-coding}), in plain
English: you cannot, on average, describe a random source using fewer
bits than its entropy; and you can get within one bit of entropy if
you are clever.  Entropy is the floor; clever coding reaches it.

Why a decision-maker should care: every storage bill, every bandwidth
bill, every ``our model is more efficient'' vendor claim bottoms out
against this floor.  If somebody tells you they compress English text
to $0.3$ bits per character, they are either wrong, lying, or have
discovered English has lower entropy than anyone thought --- which
would itself be news.  Entropy also sets the floor on forecasting:
no analyst, human or machine, can predict a genuinely high-entropy
process (next week's lottery numbers) better than chance, no matter
how much compute you throw at it.  Knowing where the entropy floor
sits for \emph{your} data is the difference between a reasonable AI
budget and a fantasy one.
\end{intuition}

Fix a finite \emph{alphabet} $\mathcal{A} = \{a_1, \dots, a_K\}$ and
let $X$ be a discrete random variable on $\mathcal{A}$ with
probability mass function $p(a) = \Pr[X = a]$.  All logarithms are
base $2$ unless stated; values are measured in \emph{bits}.

\begin{definition}[Self-information and Shannon entropy]
\label{def:entropy}
The \textbf{self-information} (or \emph{surprisal}) of the outcome
$x \in \mathcal{A}$ under $p$ is
\begin{equation}
  I_p(x) \;=\; -\log_2 p(x),
  \label{eq:self-info}
\end{equation}
with the convention $0 \cdot \log_2 0 = 0$.  The \textbf{Shannon
entropy} of $X$ is the expected self-information,
\begin{equation}
  H(X) \;=\; \mathbb{E}\bigl[I_p(X)\bigr]
       \;=\; -\sum_{k=1}^{K} p(a_k)\,\log_2 p(a_k).
  \label{eq:entropy}
\end{equation}
\end{definition}

A \textbf{code} for $\mathcal{A}$ is an injective map
$C : \mathcal{A} \to \{0,1\}^*$ assigning a binary string $C(a)$ to
each symbol.  We write $\ell(a) = |C(a)|$ for the code length.  $C$ is
a \textbf{prefix code} (or \emph{instantaneous code}) if no codeword
$C(a_i)$ is a proper prefix of any other codeword $C(a_j)$.

\begin{lemma}[Kraft's inequality]
\label{lem:kraft}
For any prefix code $C$ on $\mathcal{A}$ with lengths $\ell(a_k)$,
\begin{equation}
  \sum_{k=1}^{K} 2^{-\ell(a_k)} \;\le\; 1.
  \label{eq:kraft}
\end{equation}
Conversely, for any length assignment $\ell_1, \dots, \ell_K \in
\mathbb{N}$ satisfying~\eqref{eq:kraft} there exists a prefix code
with those lengths.
\end{lemma}

\begin{proof}[Proof sketch]
Identify each codeword with a leaf in the infinite binary tree.  A
prefix condition forbids any codeword from lying on the path to
another, so distinct codewords correspond to disjoint subtrees.  The
subtree rooted at a node of depth $\ell$ contains a $2^{-\ell}$
fraction of the infinite leaves, and disjointness gives the bound.
The converse is an explicit construction: sort $\ell_k$ in increasing
order and assign the lexicographically smallest available
length-$\ell_k$ string at each step.
\end{proof}

\begin{lemma}[Gibbs' inequality]
\label{lem:gibbs}
For any two probability mass functions $p, q$ on $\mathcal{A}$,
\begin{equation}
  -\sum_{k} p(a_k) \log_2 q(a_k)
    \;\ge\; -\sum_{k} p(a_k) \log_2 p(a_k),
  \label{eq:gibbs}
\end{equation}
with equality if and only if $p = q$.
\end{lemma}

\begin{proof}
Using $\ln x \le x - 1$ (with equality iff $x = 1$) and
$\log_2 x = \ln x / \ln 2$,
\[
  \sum_k p(a_k) \log_2 \frac{q(a_k)}{p(a_k)}
  \;\le\; \frac{1}{\ln 2}\sum_k p(a_k)\!\left(\frac{q(a_k)}{p(a_k)} - 1\right)
  \;=\; \frac{1}{\ln 2}\Bigl(\sum_k q(a_k) - 1\Bigr) \;\le\; 0,
\]
because $\sum_k q(a_k) \le 1$.  Rearranging gives~\eqref{eq:gibbs};
the equality conditions of $\ln x \le x - 1$ and $\sum q = 1$ together
force $p = q$.
\end{proof}

\begin{theorem}[Shannon source coding inequality]
\label{thm:source-coding}
Let $C$ be any uniquely decodable (in particular, any prefix) code for
$\mathcal{A}$ and let $L = \mathbb{E}[\ell(X)] = \sum_k p(a_k)\ell(a_k)$
be its expected code length.  Then
\begin{equation}
  L \;\ge\; H(X).
  \label{eq:source-coding}
\end{equation}
Moreover, there exists a prefix code (e.g.\ the Shannon--Fano code
with $\ell(a_k) = \lceil -\log_2 p(a_k)\rceil$) such that
\begin{equation}
  L \;<\; H(X) + 1.
  \label{eq:source-coding-achiev}
\end{equation}
\end{theorem}

\begin{proof}[Proof sketch]
Let $q(a_k) = 2^{-\ell(a_k)} / Z$, where
$Z = \sum_k 2^{-\ell(a_k)} \le 1$ by \cref{lem:kraft}, so $\tilde q =
q$ is a bona fide probability vector after normalisation.  By
\cref{lem:gibbs}, the cross-entropy $-\sum_k p(a_k) \log_2 \tilde
q(a_k) \ge H(X)$.  Because $-\log_2 \tilde q(a_k) = \ell(a_k) + \log_2
Z \le \ell(a_k)$ (since $Z \le 1$ makes $\log_2 Z \le 0$), we get
$L \ge -\sum_k p(a_k)\log_2 \tilde q(a_k) \ge H(X)$.  For
\eqref{eq:source-coding-achiev}, the lengths
$\lceil -\log_2 p(a_k)\rceil$ satisfy Kraft~\eqref{eq:kraft}, and each
ceiling adds at most $1$ to the expected length.  See
\citet[Ch.~5]{coverthomas2006} for the full treatment.
\end{proof}

\begin{example}[Entropy of a 4-symbol alphabet]
\label{ex:h4}
Let $\mathcal{A} = \{a, b, c, d\}$ with $p(a) = \tfrac{1}{2}$,
$p(b) = \tfrac{1}{4}$, $p(c) = p(d) = \tfrac{1}{8}$.  Then
\begin{align*}
  H(X) &= -\tfrac{1}{2}\log_2\tfrac{1}{2}
         - \tfrac{1}{4}\log_2\tfrac{1}{4}
         - 2\cdot\tfrac{1}{8}\log_2\tfrac{1}{8} \\
       &= \tfrac{1}{2}\cdot 1 + \tfrac{1}{4}\cdot 2 + \tfrac{1}{4}\cdot 3
        \;=\; \tfrac{7}{4} \text{ bits}.
\end{align*}
A prefix code achieving $L = H(X) = 1.75$ exactly is
$C(a) = 0,\ C(b) = 10,\ C(c) = 110,\ C(d) = 111$; its lengths
$(1,2,3,3)$ satisfy Kraft with equality and match
$-\log_2 p(a_k)$ symbol by symbol.
\end{example}

% =========================================================================
\section{Prediction \texorpdfstring{$\Leftrightarrow$}{<=>} Compression
         (Shannon's half of the bridge)}
\label{sec:pred-comp}
% =========================================================================

\begin{intuition}
A predictor that assigns probabilities to the next word is, mechanically,
a compressor; its error rate and its compression rate are the same number
in different clothes (see \cref{thm:ce-code} and \cref{cor:ppl}).

\textbf{The two-consultants analogy.}  Two consultants write up the
same quarter.  Both are earnest.  The one who understands your
business writes a one-pager: \emph{revenue up 4\%, driven by Zurich,
offset by FX.}  The one whose mental model is slightly off writes
ten pages of hedging --- because every surprise in the data costs
her extra sentences.  Both describe reality perfectly; only one
compresses it.  Cross-entropy is a \textbf{tax on a wrong model},
paid in pages; written $H(p,q)$ in \cref{def:cross-ent}, it is
always at least $H(p)$, and the gap is the tax.  At $0.53$ extra
bits per token over a million-token corpus that is roughly
\textbf{530{,}000 extra bits of hedging} --- a concrete, addable
number.

\textbf{Where the $0.53$ comes from.}  Suppose tomorrow's stock move
is up $90\%$ of the time, down $10\%$.  The true entropy is small,
about $0.47$ bits per day.  If your model thinks it is 50/50, your
average description length becomes $1$ bit per day.  The $0.53$-bit
gap is the per-token tax above; the million-token aggregate is the
$530{,}000$-bit number.

\textbf{Perplexity and bits per token.}  LLM papers report
``perplexity 8'' or ``$2.1$ bits per token.''  These are the same
number: $\log_2(\mathrm{PPL}) = $ bits per token = cross-entropy
(\cref{cor:ppl}).  Perplexity $8$ means the model behaves as if at
each word it were choosing uniformly among $8$ plausible options.
Lower is better.

Why a decision-maker should care: the training objective of every
major LLM \emph{is} cross-entropy on text.  When a vendor says ``our
model scored $0.1$ lower cross-entropy,'' they are saying ``our model
compresses the world $0.1$ bits better per word,'' which over
trillions of words is a lot of real modelling skill.  Two models with
identical cross-entropy are, at the level Shannon cares about, the
same model.  This gives you a vendor-agnostic yardstick --- and a
cheap test for benchmark theatre: if a new model does not improve
cross-entropy but wins on a specific test, ask whether the test was
in the training set.
\end{intuition}

Sequence models predict \emph{conditionally}.  Fix a stochastic
process $(X_t)_{t \ge 1}$ on $\mathcal{A}$ and write
$X_{<t} = (X_1, \dots, X_{t-1})$ with $X_{<1} = \emptyset$.

\begin{definition}[Cross-entropy and conditional cross-entropy]
\label{def:cross-ent}
Let $p$ be the true distribution of $(X_1, \dots, X_T)$ and let $q$ be
a model distribution (factoring as
$q(x_{1:T}) = \prod_{t=1}^{T} q(x_t \mid x_{<t})$).  The
\textbf{cross-entropy} of $q$ against $p$ over a single symbol with
history is
\begin{equation}
  H_t(p, q)
  \;=\; \mathbb{E}_{x_{<t}\sim p}\!\left[
    -\sum_{a\in\mathcal{A}} p(a \mid x_{<t})\,\log_2 q(a \mid x_{<t})
  \right],
  \label{eq:cross-ent}
\end{equation}
and the \textbf{average cross-entropy per token} is
$\overline{H}(p,q) = \tfrac{1}{T}\sum_{t=1}^{T} H_t(p,q)$.
\end{definition}

\begin{theorem}[Cross-entropy equals expected code length]
\label{thm:ce-code}
Fix $T$ and a model $q$.  Any prefix code that uses
$\ell(x_t \mid x_{<t}) = \lceil -\log_2 q(x_t \mid x_{<t})\rceil$
bits to encode $x_t$ given $x_{<t}$ satisfies
\begin{equation}
  \mathbb{E}_{x_{1:T}\sim p}\!\left[\tfrac{1}{T}\!\sum_{t=1}^{T}\ell(x_t \mid x_{<t})\right]
  \;=\; \overline{H}(p, q) + O(1/T),
  \label{eq:ce-length}
\end{equation}
and arithmetic coding attains~\eqref{eq:ce-length} with $O(1/T)$
overhead.  Moreover,
\begin{equation}
  \overline{H}(p, q)
  \;=\; \overline{H}(p) + \overline{D}_{\mathrm{KL}}(p \,\|\, q),
  \label{eq:ce-kl}
\end{equation}
where $\overline{H}(p) = \tfrac{1}{T}\sum_t H(X_t \mid X_{<t})$ is the
true conditional entropy rate and
$\overline{D}_{\mathrm{KL}}(p\|q) \ge 0$ is the average conditional
Kullback--Leibler divergence.
\end{theorem}

\begin{proof}[Proof of~\eqref{eq:ce-kl}]
For each $t$ and each history $x_{<t}$,
\[
  -\sum_{a} p(a\!\mid\! x_{<t})\log_2 q(a\!\mid\! x_{<t})
  \;=\;
  -\sum_{a} p(a\!\mid\! x_{<t})\log_2 p(a\!\mid\! x_{<t})
  + \sum_a p(a\!\mid\! x_{<t})\log_2\frac{p(a\!\mid\! x_{<t})}{q(a\!\mid\! x_{<t})}.
\]
The first term on the right is $H(X_t \mid X_{<t} = x_{<t})$; the
second is
$D_{\mathrm{KL}}(p(\cdot\mid x_{<t})\,\|\,q(\cdot\mid x_{<t}))$.
Taking expectation over $x_{<t} \sim p$ and averaging over $t$ yields
\eqref{eq:ce-kl}.  Non-negativity of the KL term is exactly
\cref{lem:gibbs}.  Statement~\eqref{eq:ce-length} follows by plugging
$\ell = \lceil -\log_2 q\rceil$ into \cref{thm:source-coding} on each
conditional distribution and absorbing the at-most-$1$-bit rounding
into the $O(1/T)$ overhead; arithmetic coding eliminates even the
integer ceilings up to a constant
\citep[Ch.~5]{coverthomas2006}.
\end{proof}

\begin{corollary}[Perplexity vs.\ bits per token]
\label{cor:ppl}
The \textbf{perplexity} of $q$ on $p$ is
$\mathrm{PPL}(q) = 2^{\overline{H}(p,q)}$, so
\begin{equation}
  \log_2 \mathrm{PPL}(q) \;=\; \overline{H}(p,q)
  \;=\; \text{(bits per token used by arithmetic coding with model }q\text{)}.
  \label{eq:ppl-bpt}
\end{equation}
\end{corollary}

\Cref{thm:ce-code} is the entire Shannon side of the bridge:
\emph{minimising the cross-entropy training objective is literally
minimising the expected compressed length}.  A next-token predictor
with cross-entropy $\overline{H}(p,q)$ bits per symbol corresponds,
via arithmetic coding, to a compressor that uses $\overline{H}(p,q)$
bits per symbol.  Better predictions $\Leftrightarrow$ shorter codes.

A worked numerical example separating the matched and mismatched
models on a 3-token sequence is given in \cref{app:ce-worked}.

% =========================================================================
\section{Kolmogorov complexity}
\label{sec:kolmogorov}
% =========================================================================

\begin{intuition}
The complexity of an object is the length of the shortest computer
program that prints it --- and nothing shorter (see
\cref{def:kolmogorov}).

\textbf{The recipe-over-the-phone analogy.}  You are a chef.  A
colleague in another city must reproduce your dish exactly.  You read
her the shortest instructions that do the job.  ``Boil water'' takes
two words.  A souffl\'{e} takes a page.  A completely random pattern
of $10{,}000$ sprinkles has no shorter description than itself: you
read out every position.  Such an object is \emph{Kolmogorov random};
it has no exploitable structure.  Everything else --- the digits of
$\pi$, the works of Shakespeare, your customer database --- has a
recipe shorter than the raw data.

\textbf{A tiny worked example.}  The string
\texttt{01010101\ldots} (a million long) has a $\sim 30$-character
recipe: \texttt{print "01" 500000 times}.  A million genuine coin
flips has no such shortcut; its recipe is roughly a million
characters.  Same disk size, wildly different Kolmogorov complexity.
Complexity measures \emph{structure}, not \emph{size}.

\textbf{Invariance and uncomputability.}  Different programming
languages differ in recipe length by at most a constant --- the size
of a translator (\cref{thm:invariance}).  So $K(x)$ is a property of
the thing, not your toolkit.  But there is no algorithm that, given
$x$, hands you its shortest program (\cref{thm:uncomputable}); if
there were, it would solve the halting problem.  $K(x)$ is a perfectly
well-defined number you can never compute exactly.

Why a decision-maker should care: $K$ is the theoretical floor every
real compressor (zip, JPEG, an LLM) approximates from above.  There is
a \emph{true} information content in your company's data, vendor-
independent.  And when we say an LLM ``understands'' a corpus, one
precise meaning is: its weights are a short program reproducing the
corpus's statistics --- the model approaching $K$ from above.
\end{intuition}

\begin{intuition}[title={The compression sniff test}]
If a vendor claims to compress genuinely random or encrypted data,
they are selling snake oil --- Shannon closed that door in 1948.  A
truly random source has entropy equal to its raw bit-length; no code
can do better on average (\cref{thm:source-coding}).  ``We compress
encrypted traffic by 40\%'' is either marketing about a non-random
side channel (timestamps, packet sizes) or a claim that should not
be in the contract.
\end{intuition}

Shannon's framework is probabilistic: entropy is a property of a
distribution, not of an individual string.  Kolmogorov's move
\citep{kolmogorov1965three} was to attach a complexity measure to a
\emph{single} object by asking for its shortest description in a fixed
universal language.

Fix once and for all a \textbf{universal Turing machine} $U$ --- a
machine that, given as input a description $\langle M, w\rangle$ of
any Turing machine $M$ and input $w$, simulates $M$ on $w$.

\begin{definition}[Kolmogorov complexity]
\label{def:kolmogorov}
The \textbf{Kolmogorov complexity} of a finite binary string
$x \in \{0,1\}^*$, relative to $U$, is
\begin{equation}
  K_U(x)
  \;=\; \min\bigl\{\, |\pi| \,:\, \pi \in \{0,1\}^*,\ U(\pi) = x \,\bigr\},
  \label{eq:kolmogorov}
\end{equation}
where $|\pi|$ is the bit-length of the program $\pi$, and $U(\pi) = x$
means $U$ on input $\pi$ halts with output $x$.
\end{definition}

\begin{theorem}[Invariance theorem]
\label{thm:invariance}
If $U_1$ and $U_2$ are two universal Turing machines, then there
exists a constant $c_{U_1,U_2}$\footnote{The constant $c_{U_1,U_2}$
is finite and does not grow with $|x|$.  For pairs of realistic
universal machines --- two reasonable programming languages
interpreting each other --- it is typically in the hundreds to low
thousands of bits.  For data objects much larger than that (a
megabyte-scale corpus, an ERP export), it is numerically negligible
and ``machine-independent up to an additive constant'' is a property
you can plan around.  For objects smaller than the constant, it is
not.}, depending on $U_1$ and $U_2$ but not on $x$, such that
\begin{equation}
  \bigl| K_{U_1}(x) - K_{U_2}(x) \bigr| \;\le\; c_{U_1,U_2}
  \quad \text{for every } x \in \{0,1\}^*.
  \label{eq:invariance}
\end{equation}
\end{theorem}

\begin{proof}[Proof sketch]
Because $U_1$ is universal, there is a fixed-length binary program
$s_{2\to 1}$ that, when prepended to any $U_2$-program $\pi$, causes
$U_1$ to simulate $U_2(\pi)$.  Thus every $U_2$-program $\pi$ of
length $|\pi|$ yields a $U_1$-program of length
$|s_{2\to 1}| + |\pi|$ with the same output, giving $K_{U_1}(x) \le
K_{U_2}(x) + |s_{2\to 1}|$.  Symmetrically with $s_{1\to 2}$.  Taking
$c_{U_1,U_2} = \max(|s_{1\to 2}|, |s_{2\to 1}|)$
gives~\eqref{eq:invariance}.
\end{proof}

By convention we henceforth write $K(x)$ and treat the universal
simulator as the additive constant.

\begin{theorem}[Uncomputability of $K$]
\label{thm:uncomputable}
There is no total computable function $f : \{0,1\}^* \to \mathbb{N}$
such that $f(x) = K(x)$ for all $x$.
\end{theorem}

\begin{proof}[Proof (Berry-paradox style)]
Suppose for contradiction such an $f$ exists.  Define the program
$P_m$, parameterised by an integer $m \in \mathbb{N}$, that does the
following: enumerate $\{0,1\}^*$ in shortlex order and, for each $x$,
compute $f(x)$ and output the first $x$ such that $f(x) > m$; then
halt.  Strings of complexity greater than $m$ exist (there are only
$2^m - 1$ programs of length less than $m$, but infinitely many
strings), so $P_m$ halts and outputs some string $x_m^\star$ with
$K(x_m^\star) > m$.

The program $P_m$ has a fixed body of bit-length $B$ (independent of
$m$) plus a binary encoding of $m$ of at most
$\lceil \log_2(m+1)\rceil + O(1)$ bits (e.g.\ Elias delta).  So
$|P_m| \le B + \log_2 m + O(1)$, giving
\begin{equation}
  K(x_m^\star) \;\le\; |P_m| \;\le\; B + \log_2 m + O(1).
  \label{eq:berry-upper}
\end{equation}
Combined with $K(x_m^\star) > m$,
\begin{equation}
  m \;<\; B + \log_2 m + O(1),
  \label{eq:berry-contra}
\end{equation}
which fails for every sufficiently large $m$ (since $m - \log_2 m \to
\infty$), a contradiction.
\end{proof}

The upshot of \cref{thm:uncomputable} is that every practical
compressor --- gzip, PNG, FLAC, and indeed every neural language
model --- produces an upper bound on $K(x)$, never its exact value.

% =========================================================================
\section{Solomonoff induction}
\label{sec:solomonoff}
% =========================================================================

\begin{intuition}
To predict the future, average the predictions of every possible
theory, giving shorter theories exponentially more weight
(\cref{def:M,thm:solomonoff}).

\textbf{The prediction-market analogy.}  Imagine a prediction market
where every trader is a theory of the world.  Simple theories start
with a large float; complex theories start with almost nothing ---
their share price is their \emph{simplicity premium}, set by theory
length: a theory of length $L$ bits gets weight $2^{-L}$.  A theory
you can state in 10 words gets a float $1{,}000$ times larger than
one needing 20 words.  Each new data point is an earnings call.
Theories that predicted it cheaply get paid; theories that predicted
it badly get margin-called.  Over time the market consolidates around
the simplest theories that still explain the tape.  That
consolidation \emph{is} the Solomonoff prediction $M$.  This is
Occam's razor, precisely quantified.

\textbf{A tiny worked example.}  You see $2, 4, 6, 8$.  What's next?
The theorist ``add 2'' has a very short program and shouts
``$10$'' at high volume.  The theorist ``the sequence is
$2,4,6,8,47$, then random'' has a longer program and shouts ``$47$''
much more quietly.  Weighted vote: overwhelmingly ``$10$.''  If the
next number really is $47$, all the ``add 2'' theorists are silenced
forever, and the surviving theorists with the shortest explanation of
$2,4,6,8,47$ take over.  The machine self-corrects as evidence
arrives.  A toy two-hypothesis universe is worked out in
\cref{ex:toy-solomonoff}.

\textbf{Convergence.}  If the true process is computable, Solomonoff's
predictor converges to it: KL error goes to zero
(\cref{thm:solomonoff}).  The intuition: the true theory is somewhere
in the panel, and every wrong theorist loses microphone share each
time she makes a wrong prediction.  Eventually the truth dominates.

Why a decision-maker should care: Solomonoff induction is the gold
standard of ``learning from data.''  Every ML method is, badly and
efficiently, trying to imitate it.  ``Bayesian model averaging,''
``ensemble methods,'' ``regularise for simplicity'' --- all are
engineering approximations of this ideal.  It also explains why
overfitting is bad in a principled way: an overfit model is a long,
ad-hoc theory, and Solomonoff would have weighted it near-zero from
the start.  The question ``which model is
best?'' has a right answer: whichever predicts held-out data with the
fewest bits of surprise.
\end{intuition}

Kolmogorov complexity selects a single shortest program.  Solomonoff's
prior \citep{solomonoff1964formal} performs a \emph{Bayesian average}
over all programs, weighted by their lengths.  Fix a
\emph{prefix-free} universal Turing machine $U$: the set of halting
inputs forms a prefix code, so
$\sum_{\pi : U(\pi)\downarrow} 2^{-|\pi|} \le 1$ by Kraft's inequality.

\begin{definition}[Universal prior $M$]
\label{def:M}
For a finite binary string $x \in \{0,1\}^*$, the \textbf{universal
(Solomonoff) prior} is
\begin{equation}
  M(x) \;=\; \sum_{\pi \,:\, U(\pi)\, =\, x *} 2^{-|\pi|},
  \label{eq:M}
\end{equation}
where the sum runs over all programs $\pi$ such that $U(\pi)$ halts
and outputs a string whose prefix is $x$ (denoted $x*$).
Equivalently, $M(x)$ is the probability that $U$ fed with a uniformly
random infinite bitstream prints $x$ as its first $|x|$ output bits.
\end{definition}

The Solomonoff predictive posterior is
\begin{equation}
  M(x_{t+1} \mid x_{1:t}) \;=\; \frac{M(x_{1:t}\,x_{t+1})}{M(x_{1:t})},
  \label{eq:M-posterior}
\end{equation}
defined whenever $M(x_{1:t}) > 0$.

\begin{theorem}[Solomonoff's convergence theorem]
\label{thm:solomonoff}
Let $\mu$ be any computable probability measure on $\{0,1\}^\infty$.
Then there is a constant
\begin{equation}
  c_\mu \;=\; 2^{-K(\mu) + O(1)}
  \label{eq:cmu}
\end{equation}
such that, for all finite prefixes $x$,
\begin{equation}
  M(x) \;\ge\; c_\mu \cdot \mu(x).
  \label{eq:dominance}
\end{equation}
Consequently, if $X_{1:\infty} \sim \mu$, the cumulative expected
KL divergence between the Solomonoff predictor and $\mu$ is bounded
uniformly:
\begin{equation}
  \sum_{t=1}^{\infty}
    \mathbb{E}_{\mu}\!\left[
      D_{\mathrm{KL}}\!\bigl(\mu(\cdot \mid X_{<t}) \,\|\, M(\cdot \mid X_{<t})\bigr)
    \right]
  \;\le\; K(\mu)\ln 2 + O(1).
  \label{eq:sol-bound}
\end{equation}
In particular,
$D_{\mathrm{KL}}(\mu(\cdot \mid X_{<t}) \| M(\cdot \mid X_{<t})) \to 0$
in $\mu$-expectation as $t \to \infty$.
\end{theorem}

\begin{proof}[Proof sketch]
Because $\mu$ is computable, it has a description as a program
$\pi_\mu$ of length $|\pi_\mu| = K(\mu) + O(1)$.  Fix any family of
programs that, given $\pi_\mu$ as subroutine, simulates a $\mu$-sample
to within arbitrarily many bits; summing $2^{-|\pi|}$ over that
family and restricting to those whose output begins with $x$ shows
$M(x) \ge 2^{-K(\mu) - O(1)}\,\mu(x)$, which is
\eqref{eq:dominance} with $c_\mu$ as in~\eqref{eq:cmu}.  Taking
logarithms and applying a telescoping chain-rule argument to the KL
divergence then yields~\eqref{eq:sol-bound}.  The correct program
$\pi_\mu$ has prior mass exactly $2^{-K(\mu)}$; the finite KL budget
is the information-theoretic price of identifying it.  For the full
proof see \citet[Ch.~5]{livitanyi2019}.
\end{proof}

\begin{example}[Toy Solomonoff update on a two-program universe]
\label{ex:toy-solomonoff}
Restrict the universe of programs to two deterministic generators:
$\pi_0 : \text{always output }0$ (bit-length $|\pi_0| = 2$) and
$\pi_1 : \text{output the cycle } 01010101\ldots$ (bit-length
$|\pi_1| = 3$).  Under the Solomonoff-style weighting
$w(\pi) = 2^{-|\pi|}$, the prior over the two hypotheses is
\[
  P(\pi_0) \;=\; \frac{2^{-2}}{2^{-2} + 2^{-3}} \;=\; \frac{2}{3},
  \qquad
  P(\pi_1) \;=\; \frac{2^{-3}}{2^{-2} + 2^{-3}} \;=\; \frac{1}{3}.
\]
Observe the prefix $x_{1:4} = 0101$.  Likelihoods:
$\Pr(0101 \mid \pi_0) = 0$ (the all-zeros generator cannot produce a
$1$); $\Pr(0101 \mid \pi_1) = 1$.  Bayes' rule gives
$P(\pi_0 \mid 0101) = 0$, $P(\pi_1 \mid 0101) = 1$.  So the Solomonoff
predictor assigns $M(x_5 = 0 \mid 0101) = 1$, $M(x_5 = 1 \mid 0101) =
0$.  Two features: the shorter program $\pi_0$ was \emph{not}
favoured after data --- simplicity is a prior preference overridden
by evidence --- and a single contradicting bit collapses an
exponentially-weighted prior to a point mass, which is the mechanism
by which Solomonoff induction converges so fast on deterministic
sources.
\end{example}

% =========================================================================
\section{AIXI: intelligence as optimal compression-driven prediction}
\label{sec:aixi}
% =========================================================================

\begin{intuition}
The theoretically optimal agent is the one that, before each action,
asks the Solomonoff panel what is likely to happen, and picks the
action whose expected future reward is highest
(\cref{def:aixi,thm:aixi-opt}).

\textbf{The regional-manager analogy.}  Imagine a regional manager
who, before every staffing decision, consulted every possible theory
of her stores --- every demand model, every weather-effect hypothesis,
every crank manifesto about Mercury retrograde --- weighted by how
simple each theory is.  For each candidate decision she asks: ``If I
do this, what happens over the next hundred Tuesdays, and how much
will I enjoy it?''  She takes the weighted average and places the
staffing bet that maximises expected sales across all of them.  That
manager would be provably optimal in a precise technical sense.  She
would also be impossible: consulting every possible theory takes
infinite time.  This is \emph{AIXI}.  It is a \textbf{north star for
what a learner is approximating, not a machine anyone can build}.
Real systems are food trucks reading the Michelin menu; AIXI is the
chef with an infinite pantry and no kitchen.

\textbf{A tiny worked example.}  A two-armed bandit: Arm A paid out
once in three pulls; Arm B paid out zero times in one pull.  The
simplest theories favour A; slightly more complex theories say B was
just unlucky.  Weight theories by simplicity, compute expected reward
from each arm over the next $100$ pulls, pick the winner.  You will
usually pull A but occasionally explore B --- not because someone
programmed exploration, but because the simplicity-weighted average
\emph{naturally} keeps B alive.  AIXI does this for every decision,
everywhere.

\textbf{Why Hutter calls it ``universal intelligence.''}  AIXI does
not need to know the game, the physics, or the rules.  Drop it into
any computable environment and it will eventually learn the
environment and play optimally.  Its only prior is Occam's razor.

\textbf{The catch.}  AIXI is uncomputable; it requires running every
possible program.  Bounded approximations exist but are toys at
present.  AIXI is to intelligence what the ideal gas law is to a
real gas --- a north star, not a product.

Why a decision-maker should care: AIXI is a crisp, non-hand-wavy
definition of ``what would a perfectly general AI do?''  When someone
claims their system is ``a step toward AGI,'' ask: is it a step
toward better approximating AIXI?  Is it building better world-models
and planning against them?  If yes, the claim is in the right
ballpark.  If it is pattern-matching in a narrow domain with no
planning layer, the claim is marketing.  And no real system \emph{is}
AIXI, so any product advertised as ``universal intelligence'' is
abusing the term.
\end{intuition}

Hutter's AIXI framework \citep{hutter2005universal} extends Solomonoff
induction from passive sequence prediction to sequential
decision-making.  The agent interacts with an environment over
discrete time steps.  At step $t$ the agent chooses an action $a_t
\in \mathcal{A}_{\mathrm{act}}$, and the environment returns a
percept $e_t = (o_t, r_t)$ consisting of an observation $o_t \in
\mathcal{O}$ and a reward $r_t \in [0, r_{\max}] \subset \mathbb{R}$.
Write the interaction history as
$h_{<t} = (a_1, e_1, a_2, e_2, \dots, a_{t-1}, e_{t-1})$.

An \textbf{environment} is a conditional probability measure
$\nu(e_{1:m} \mid a_{1:m})$ over percept sequences given action
sequences.  The \textbf{universal mixture environment} is
\begin{equation}
  \xi(e_{1:m} \mid a_{1:m})
  \;=\; \sum_{\nu \in \mathcal{M}} 2^{-K(\nu)}\,\nu(e_{1:m} \mid a_{1:m}),
  \label{eq:xi}
\end{equation}
where $\mathcal{M}$ is the class of all lower-semicomputable
(``chronological'') environments and $K(\nu)$ is the Kolmogorov
complexity of a prefix program for $\nu$.

\begin{definition}[AIXI]
\label{def:aixi}
Fix a finite planning horizon $m$ and a discount.  The \textbf{AIXI
action} at history $h_{<t}$ is
\begin{equation}
  a_t^{\text{AIXI}}
  \;=\; \arg\max_{a_t}
        \sum_{e_t}\!\max_{a_{t+1}}\sum_{e_{t+1}}\!\cdots\!\max_{a_m}\sum_{e_m}
        \Bigl(\sum_{k=t}^{m} r_k\Bigr)\,
        \xi\bigl(e_{t:m} \,\big|\, h_{<t}, a_{t:m}\bigr).
  \label{eq:aixi-rule}
\end{equation}
\end{definition}

In words: AIXI plans by expectimax over all possible futures, weighted
by the Solomonoff-type universal mixture~\eqref{eq:xi}.  The
$2^{-K(\nu)}$ weighting is what ties the definition of intelligence
back to compression: hypothesised environments with short descriptions
receive exponentially more weight, and an AIXI-optimal agent
simultaneously (i) compresses its environment optimally and (ii) acts
to maximise reward under that compressed model.

\begin{theorem}[Universal optimality of AIXI, informal]
\label{thm:aixi-opt}
Within the class of lower-semicomputable environments, AIXI is
\emph{Pareto optimal}: no agent $\pi'$ achieves strictly greater
expected total reward than AIXI in every environment $\nu \in
\mathcal{M}$.  Moreover, for any computable $\nu$ the AIXI predictor
converges to the true environment's predictions in the sense of
\eqref{eq:sol-bound} with $M$ replaced by~$\xi$.
\end{theorem}

For the precise statement and proof see
\citet[Ch.~5]{hutter2005universal}.  The significance of
\cref{thm:aixi-opt} is that ``intelligence'' has, within this
framework, been reduced to a single ingredient (the universal
prior~\eqref{eq:xi}) plus bookkeeping (the expectimax expansion).
Remove the Solomonoff prior and the agent collapses to uninformed
search; keep it, and the agent becomes universally optimal.

% =========================================================================
\section{What this buys us for LLMs}
\label{sec:llms}
% =========================================================================

\begin{intuition}
An LLM trained to predict the next word is, mechanically and provably,
a compressor of its training data --- which, per the chain above,
puts it somewhere on the intelligence spectrum
(\cref{eq:chain,thm:ce-code,thm:solomonoff}).

\textbf{Residuals and arithmetic coding, in plain English.}  A
\emph{residual} is the gap between what the model predicted and what
actually came next.  If the model gave next-token probability $0.9$
to ``the'' and ``the'' showed up, the residual is tiny --- a fraction
of a bit.  If the model gave $0.01$ and ``the'' still showed up, the
residual is large --- about $6.6$ bits of surprise.
\emph{Arithmetic coding} is just a near-optimal way to write all
those surprise numbers down end to end.  The compressed file is
therefore the model's weights (fixed cost) plus the bill for every
surprise (variable cost).  ``KL divergence'', when it appears below,
is the same tax formalised: the average extra bits per token you pay
for using the wrong distribution.

\textbf{The Hutter Prize as a concrete hinge.}  Marcus Hutter put up
real prize money for whoever can compress 1~GB of Wikipedia the most.
The current leaderboard is dominated by neural compressors.  Training
an LLM on Wikipedia to minimise cross-entropy \emph{is} compressing
Wikipedia; the model weights plus arithmetic-coded residuals are the
compressed file.  Better prediction $=$ smaller file $=$ (by
\cref{thm:ce-code}) better distributional model $=$ (by
\cref{thm:solomonoff}) closer to the true generating process $=$
(per Hutter) more intelligence.

\textbf{Scaling laws.}  Plot model size against cross-entropy loss on
held-out text: you get a strikingly clean curve.  Double the data,
double the parameters, loss drops predictably.  The curve has been
stable across four orders of magnitude
\citep{kaplan2020scaling}.  Cross-entropy --- our bits-per-token
yardstick --- is a smooth function of investment.  You can buy
intelligence by the bit.  Not efficiently, not forever, and with
diminishing returns --- but predictably.

\textbf{A tiny worked example.}  GPT-2 hit roughly $1.2$ bits per
English character on a standard benchmark; GPT-4-class models are
closer to $0.7$.  That looks like a modest drop.  But English text is
about $1.3$ bits of true entropy per character (Shannon's 1950s
estimate).  So GPT-2 was $\sim 92\%$ of the way to the floor; GPT-4
is $\sim 98\%$.  Small in bits, enormous in dollars --- the last bit
contains the reasoning, the rare facts, the nuanced world-models.

\textbf{Emergent capabilities, cautious framing.}  If the training
objective is a compression proxy and compression is provably linked
to a universal notion of intelligence, then LLMs doing arithmetic,
translation, and code --- capabilities nobody explicitly trained for
--- is not a bug.  It is the theory predicting itself.  To compress
text well, you must model the world the text is about.

Why a decision-maker should care, three upshots.  First, you now have
a real yardstick --- cross-entropy on your own domain data --- to
compare vendors.  Second, you understand why ``just predict the next
word'' produces general-looking behaviour; it is not marketing, the
mathematics agrees.  Third, you can reason about the
ceiling: as models approach the true entropy of text, progress on
pure language tasks must slow, and the frontier moves to multimodal
data, embodied interaction, and reasoning architectures.  Plan
capital allocation accordingly.
\end{intuition}

Concretely, a large language model parameterises a conditional
distribution $q_\theta(x_t \mid x_{<t})$ and is trained by minimising
the negative log-likelihood
\begin{equation}
  \mathcal{L}(\theta)
  \;=\; -\frac{1}{T}\sum_{t=1}^{T} \log_2 q_\theta(x_t \mid x_{<t}).
  \label{eq:llm-loss}
\end{equation}
By \cref{thm:ce-code}, $\mathcal{L}(\theta)$ is the expected
arithmetic-coding length per token: minimising cross-entropy \emph{is}
minimising compressed size.  Equivalently, by \cref{cor:ppl}, lowering
perplexity lowers bits-per-token by the same amount.

\citet{kaplan2020scaling} showed that $\mathcal{L}(\theta)$ follows
clean power laws in model size $N$, dataset size $D$, and compute $C$,
of the form $\mathcal{L} \approx (N_c / N)^{\alpha_N}$.  Viewed
through~\eqref{eq:llm-loss}, each such power-law decrement is a
\emph{bit-per-token compression improvement}.  The operational
content of \cref{thm:solomonoff} in this setting is that, to the
extent that the data-generating process is computable, any predictor
whose bits-per-token approaches the conditional entropy rate
$\overline{H}(p)$ is implicitly approximating the $\mu$-posterior of
the universal prior $M$.

\citet{deletang2024language} operationalised this equivalence by
wiring a large language model into an arithmetic coder and
compressing raw bytes of image and audio data.  They reported
compression ratios lower than domain-specific specialist compressors
(e.g.\ PNG, FLAC) on out-of-domain modalities, despite the model
having been trained solely on text.  From the vantage point of
\cref{thm:ce-code,thm:solomonoff}, the experiment is a consistency
check: a predictor that approximates a universal prior should, by
\eqref{eq:dominance}, compress any computable source to within an
additive $K(\mu)$ of its true entropy, and text exposure is enough
to learn a useful approximation of that prior.

The \textbf{Hutter Prize} \citep{hutter2005universal} formalises the
economic version of the claim: a cash prize for improving the
lossless compression ratio on a fixed text corpus.  Within the
equivalence framework, every winning submission is, by definition, a
model whose cross-entropy on the corpus is lower --- which, by
\cref{cor:ppl}, is a more accurate next-symbol predictor.

Collecting the theorems of
\crefrange{sec:pred-comp}{sec:aixi} we obtain the \emph{working
equivalence}:
\begin{equation}
  \underbrace{\text{low cross-entropy}}_{\text{good predictor}}
  \;\Longleftrightarrow\;
  \underbrace{\text{short arithmetic code}}_{\text{good compressor}}
  \;\Longleftrightarrow\;
  \underbrace{\text{approx.\ of universal prior}}_{\substack{\text{Solomonoff/AIXI}\\\text{sense of intelligence}}}
  \label{eq:chain}
\end{equation}
The double arrows are theorems within Shannon's framework
(\cref{thm:ce-code}) and asymptotic theorems modulo computability
(\cref{thm:solomonoff}).

% =========================================================================
\section{Caveats and what the equivalence does not say}
\label{sec:caveats}
% =========================================================================

\begin{intuition}
The equivalence is real but narrow: it tells you that good prediction
entails good \emph{modelling}, not that good prediction entails
understanding, agency, or safety.  It promises statistical competence.
It does not promise causal reasoning, counterfactual simulation, or
moral judgement.  Those may emerge --- the theory leaves the door
open --- but they are not guaranteed by the objective.

\textbf{Uncomputability is not cosmetic.}  $K(x)$ and AIXI are
unreachable (\cref{thm:uncomputable}).  Every real compressor and
every real agent is a finite approximation.  Saying an LLM is
``approximating AIXI'' is like a food-truck claiming to approximate
a Michelin three-star kitchen by reading the menu --- technically
on the same axis, in practice not the relevant comparison.

\textbf{The only non-tautological falsifier.}  A reasonable operator
will ask: if AIXI is uncomputable and every deployed system is ``a
finite approximation'', what test other than cross-entropy itself
decides that system A approximates the ideal better than system B?
Honest answer: held-out cross-entropy on distributions the training
set did not see, plus downstream task transfer.  There is no deeper
falsifier in the current theory.  If that makes the AIXI framing
feel tautological to an operator, the operator is reading the state
of the field correctly.

\textbf{Compression $\ne$ understanding.}  A junior analyst who has
memorised every deal memo in the vault --- say, $4{,}000$ deals
going back a decade --- has compressed a lot of data.  Ask her to
price a deal that is not in the vault and you find out whether she
compressed the memos or the market.  A lossless compressor has
captured statistics; that is \emph{some} kind of understanding, but
not necessarily the kind a manager means when she asks ``does the
model understand our business?''

\textbf{Embodiment.}  A text-only system has only ever been surprised
by text.  It has no ground-truth feedback from physical reality.
Solomonoff's framework accepts any computable environment, but LLMs
have the text channel and lack the body-in-a-world channel.  Agentic,
tool-using, multimodal systems are early attempts to close that gap.

\textbf{Benchmarks are not the floor.}  A model with low cross-entropy
on a benchmark may just have memorised the benchmark.  Data
contamination is rampant.  The equivalence only holds when the test
distribution is genuinely held out.

Why a decision-maker should care: use the equivalence as a
\emph{lower bound}, not an upper bound.  When a vendor says ``lowest
perplexity,'' accept that as real evidence of genuine modelling skill
--- and then ask the second question: skill at what, tested on what,
does it translate to your domain?  The equivalence tells you the knob
exists.  It does not tell you that turning the knob solves your
problem.  Confusing ``next-token prediction'' with ``just autocomplete''
misses a deep result; confusing ``next-token prediction'' with
``general intelligence'' overclaims a narrow one.  Both errors are
expensive.  Staying in the middle --- taking the equivalence seriously
but not theologically --- is the managerial win.
\end{intuition}

Several limitations of the chain~\eqref{eq:chain} deserve formal
acknowledgement.

\paragraph{(C1) Uncomputability.}
By \cref{thm:uncomputable}, $K(x)$ is not computable; by a related
argument so is $M(x)$.  The entire right-hand side of~\eqref{eq:chain}
therefore names a mathematical ideal, not an algorithm.  Every real
compressor produces an upper bound $\tilde K(x) \ge K(x)$, and
equality is in principle unattainable.  ``Approximates the universal
prior'' must be taken as an asymptotic statement about cross-entropy,
not as a literal claim that the model is computing~$M$.

\paragraph{(C2) The additive constant.}
\Cref{thm:invariance} only guarantees machine-independence up to an
additive constant $c_{U_1, U_2}$.  For short strings this constant
can dominate $K(x)$ itself; Kolmogorov complexity is only meaningful
asymptotically.  Analogously, the $K(\mu)$ term in~\eqref{eq:sol-bound}
is a finite one-time information cost: it does not shrink with data.

\paragraph{(C3) Practical compressors vs.\ $K$.}
A compressor $C$ is \emph{universal} in Shannon's sense if its code
length $L_C(x_{1:T})/T$ converges to the entropy rate of any
stationary ergodic source.  It is \emph{not} in general universal in
Solomonoff's sense: it need not dominate all computable measures.  The
gap between Lempel--Ziv (Shannon-universal) and Solomonoff (universal
over all computable measures) is a full logical level of universality,
and neural compressors live somewhere in between.

\paragraph{(C4) Necessary vs.\ sufficient.}
\Cref{thm:solomonoff,thm:aixi-opt} establish that optimal prediction
and agency \emph{entail} optimal compression.  They do \emph{not}
establish that every improvement in compression corresponds to an
improvement in every intuitive sense of ``understanding.''  In
particular, the framework is silent on phenomenal consciousness,
grounding, and embodiment, and it takes the reward signal and the
action--percept loop as primitives rather than explaining their
origin.

\paragraph{(C5) Compression $\ne$ understanding (as internal structure).}
A compressor is characterised by its length function, not by its
internal representations.  Two compressors with identical code
lengths may use entirely different internal mechanisms, only one of
which we would intuitively call ``understanding.''
Equation~\eqref{eq:chain} reduces ``intelligence'' to an extensional
(input--output) property; any argument that understanding is
intensional (about how the function is computed) is not touched by
the chain.

\paragraph{Summary.}
The equivalence~\eqref{eq:chain} is a genuine theorem modulo
\crefrange{def:entropy}{def:aixi} and the computability-based caveats
(C1)--(C3).  What it licences is a precise sense in which lowering
cross-entropy on a sufficiently rich corpus moves a predictor along a
single axis shared with compression and with Hutter-style universal
agency.  What it does \emph{not} licence is the stronger claim that
this axis exhausts what we mean by understanding.

% =========================================================================
\section*{Further reading}
\addcontentsline{toc}{section}{Further reading}
% =========================================================================

\begin{description}[leftmargin=1em,style=nextline]
  \item[Information theory, foundations.] \citet{shannon1948mathematical}
    is the founding paper and still the clearest exposition of the
    source-coding setup.  \citet{coverthomas2006} is the modern textbook
    reference; chapters 2--5 cover everything in
    \cref{sec:shannon-prelude,sec:pred-comp}.
  \item[Kolmogorov complexity and Solomonoff induction.]
    \citet{kolmogorov1965three} introduces the three approaches to
    defining information; \citet{solomonoff1964formal} introduces the
    universal prior.  \citet{livitanyi2019} is the definitive modern
    textbook, with full proofs of \cref{thm:invariance,thm:uncomputable,thm:solomonoff}.
  \item[Universal agency.] \citet{hutter2005universal} is the original
    AIXI monograph, including proofs of \cref{thm:aixi-opt} and
    related optimality results.
  \item[LLMs as compressors, empirical.]
    \citet{kaplan2020scaling} gives the scaling-law curves discussed
    in \cref{sec:llms}.  \citet{deletang2024language} demonstrates
    experimentally that a text-trained large language model, used as
    the model inside an arithmetic coder, compresses out-of-domain
    byte streams (images, audio) competitively.
\end{description}

\appendix

% =========================================================================
\section{Worked cross-entropy calculation}
\label{app:ce-worked}
% =========================================================================

\begin{example}[Cross-entropy on a 3-token sequence; nats vs.\ bits]
\label{ex:ce3}
Let $\mathcal{A} = \{A,B,C\}$ and suppose the true conditional
distributions at positions $t=1,2,3$ have already been conditioned on
the observed prefix, giving the sequence of true posterior masses on
the realised tokens $p_1 = 0.8,\ p_2 = 0.5,\ p_3 = 0.1$.  Consider
two models.  Model $q^{(1)}$ assigns $q_1 = 0.8, q_2 = 0.5,
q_3 = 0.1$ (matches truth on the realised tokens); model $q^{(2)}$
assigns $q_1 = 0.1, q_2 = 0.1, q_3 = 0.8$ (systematically wrong).

For model $q^{(1)}$, the per-token code lengths in bits are
$-\log_2 0.8 \approx 0.3219$, $-\log_2 0.5 = 1$,
$-\log_2 0.1 \approx 3.3219$, summing to approximately $4.644$~bits.
In nats, multiply by $\ln 2 \approx 0.6931$: $4.644 \cdot \ln 2
\approx 3.219$~nats, matching $-\ln(0.8 \cdot 0.5 \cdot 0.1) =
-\ln 0.04 \approx 3.219$ directly.

For model $q^{(2)}$, the lengths are $-\log_2 0.1 \approx 3.3219$,
$-\log_2 0.1 \approx 3.3219$, $-\log_2 0.8 \approx 0.3219$, totalling
$\approx 6.966$~bits ($\approx 4.828$~nats, consistent with
$-\ln(0.008) \approx 4.828$).  The wrong model spends $2.32$ extra
bits, the Kullback--Leibler excess predicted by~\eqref{eq:ce-kl}.
Conversion rule: $1\text{ nat} = \log_2 e \text{ bits} \approx
1.4427$~bits and conversely $1\text{ bit} = \ln 2 \text{ nats}$.
\end{example}

\bibliography{../section3}

\end{document}
